% Ch3 WS1 -- moles, molar mass, percent composition, formulas. Blocks 1-2.
\documentclass{apchem-worksheet}

\apcunitword{Chapter}
\apcunit{3}
\apcunittitle{Stoichiometry}
\apcdoctitle{Worksheet 1 \textbullet\ Moles and Formulas}
\apcsubtitle{Zumdahl \S3.1--3.7 \textbullet\ every answer needs units and correct sig figs}

\begin{document}
\wsheader[Show the dimensional-analysis setup, not just the number. An
answer with no units is wrong even when the digits are right.]

\problem[]{Molar masses --- build each from the periodic table:}
\begin{parts}
  \item \ce{Mg(OH)2}: \answerblank[3.4cm]{\SI{58.33}{\gram\per\mole}}
  \item \ce{Al2(SO4)3}: \answerblank[3.4cm]{\SI{342.17}{\gram\per\mole}}
  \item \ce{C6H12O6}: \answerblank[3.4cm]{\SI{180.16}{\gram\per\mole}}
\end{parts}

\problem[]{Conversions:}
\begin{parts}
  \item Moles in \SI{54.0}{\gram} \ce{H2O}:
        \answerblank[3cm]{\SI{3.00}{\mole}}
  \item Mass of \SI{0.250}{\mole} \ce{CaCO3}:
        \answerblank[3cm]{\SI{25.0}{\gram}}
  \item Molecules in \SI{11.0}{\gram} \ce{CO2}:
        \answerblank[3cm]{$1.50\times10^{23}$}
  \item Oxygen \emph{atoms} in \SI{0.100}{\mole} \ce{Al2(SO4)3}:
        \answerblank[3cm]{$7.23\times10^{23}$}
\end{parts}

\problem[]{A sample of an unknown element has a mass of \SI{4.00}{\gram} and
contains $9.03\times10^{22}$ atoms.}
\begin{parts}
  \item Moles present: \answerblank[3cm]{\SI{0.150}{\mole}}
  \item Molar mass: \answerblank[3cm]{\SI{26.7}{\gram\per\mole}}
  \item Most likely identity: \answerblank[3.7cm]{aluminum (26.98)}
\end{parts}

\problem[]{Percent composition:}
\begin{parts}
  \item $\%\,\ce{N}$ in \ce{NH4NO3}: \answerblank[2.8cm]{35.00\%}
  \item $\%\,\ce{O}$ in \ce{Al2(SO4)3}: \answerblank[2.8cm]{56.11\%}
  \item Mass of nitrogen in \SI{50.0}{\gram} of \ce{NH4NO3}:
        \answerblank[2.8cm]{\SI{17.5}{\gram}}
\end{parts}

\problem[]{A compound is 26.6\% K, 35.4\% Cr, and 38.1\% O by mass.}
\begin{parts}
  \item Empirical formula: \workspace[2.8cm]
        \keyonly{\begin{apcanswer}K: $26.6/39.10 = 0.680$;
        Cr: $35.4/52.00 = 0.681$; O: $38.1/16.00 = 2.38$.
        Divide by 0.680: $1 : 1 : 3.5 \to \times2 \to$
        \ce{K2Cr2O7}\end{apcanswer}}
  \item Why must the 3.5 be doubled rather than rounded to 4?
        \answerlines[2]{A ratio ending in .5 is a signal that the true
        subscripts are twice these values; rounding would change the
        compound's composition.}
\end{parts}

\problem[]{A hydrocarbon is 82.7\% C and 17.3\% H by mass, with a molar mass
of \SI{58.12}{\gram\per\mole}.}
\begin{parts}
  \item Empirical formula: \answerblank[2.6cm]{\ce{C2H5}}
  \item Molecular formula: \answerblank[2.6cm]{\ce{C4H10}}
\end{parts}

\problem[]{Combustion analysis. A \SI{0.2000}{\gram} sample of a compound
containing only C, H, and O yields \SI{0.4547}{\gram} \ce{CO2} and
\SI{0.1862}{\gram} \ce{H2O}. Carry four significant figures throughout ---
combustion data is unforgiving of early rounding.}
\begin{parts}
  \item Mass of C: \answerblank[3cm]{\SI{0.1241}{\gram}}
  \item Mass of H: \answerblank[3cm]{\SI{0.02083}{\gram}}
  \item Mass of O (by difference): \answerblank[3cm]{\SI{0.0551}{\gram}}
  \item Empirical formula: \workspace[2.6cm]
        \keyonly{\begin{apcanswer}C: $0.01033$; H: $0.02066$; O: $0.003444$
        mol. Divide by 0.003444: C 3.00, H 6.00, O 1.00 $\to$
        \ce{C3H6O}.\end{apcanswer}}
  \item If the molar mass is \SI{58.08}{\gram\per\mole}, give the molecular
        formula. \answerblank[3cm]{\ce{C3H6O} (same)}
\end{parts}

\begin{spiralstrip}{Chapter 2 \textbullet\ keep naming warm}
  \item Name \ce{Fe2(SO4)3}: \answerblank[4.5cm]{iron(III) sulfate}
  \item Formula for ammonium carbonate:
        \answerblank[3cm]{\ce{(NH4)2CO3}}
  \item Name \ce{N2O4}: \answerblank[4.5cm]{dinitrogen tetroxide}
  \item Electrons in \ce{Cr^3+}: \answerblank[2cm]{21}
  \item Which law: \ce{CO} and \ce{CO2} oxygen ratio is 1:2?
        \answerblank[4.5cm]{multiple proportions}
\end{spiralstrip}

\end{document}
